\subsection{The source torus-to-proper-group constancy argument}
\label{subsec:faithful-torus-proper-constant}

The quotient-action producer uses the following result from
\texttt{SplitTorusToAbelianVarietyConst}.  We give its actual
two-chart and resultant argument.

\begin{lemma}\label{lem:faithful-torus-proper-constant}
Let \(K\) be a field and \(B\) a group scheme proper over \(K\).
Let \(X\) be a finite free abelian group.  Every \(K\)-morphism
\[
                         D(X)\longrightarrow B
\]
is constant, with value its value at the identity of the split
torus \(D(X)\).  In particular a group homomorphism is zero.
The map in the first assertion need not be a homomorphism.
\end{lemma}
\begin{proof}
We first treat one Laurent variable over any field extension
\(L/K\).  A morphism
\(\operatorname{Spec}L[T,T^{-1}]\to B\) extends over
\(\operatorname{Spec}L[T]\): every local ring of the latter
regular integral curve is a valuation ring, and the valuative
criterion for properness extends the map from the function
field at each point.  Separatedness makes the extensions agree
on overlaps.  Apply the same construction after \(T\mapsto T^{-1}\).
We obtain two maps
\[
        b,b':\mathbb A^1_L\longrightarrow B,\qquad
                         b(x)=b'(x^{-1})\quad(x\ne0),
 \tag{TP1}\label{eq:faithful-torus-two-charts}
\]
with this equality after arbitrary field extension.

Here is the source's algebraic neighbourhood calculation.
Let \(R\) be a \(K\)-algebra with a surjective specialization
\(e:R\to L\).
Suppose \(F,G:\mathbb A^1_R\to B\) satisfy the analogue of
\eqref{eq:faithful-torus-two-charts} on all field-valued points,
and both specializations have image in an affine open
\(V\subset B\).
The closed complement of \(F^{-1}V\) is cut out by an ideal of
\(R[T]\) whose image in \(L[T]\) is the unit ideal.
Surjectivity of \(R[T]\to L[T]\) therefore supplies an element
\(a\) of that ideal with \(e(a)=1\).
In particular, on any field-valued point where \(F(x)\notin V\),
one has \(a(x)=0\).
Likewise choose \(a'\in R[T]\) with \(e(a')=1\) vanishing where
\(G\) leaves \(V\).

Put \(n=\deg a'\) and
\[
                  h(T)=T^{n+1}a'(T^{-1}).
\]
This is a polynomial, with \(h(0)=0\).
Its resultant with \(a\), taken using the indicated degree
bounds, gives \(g\in R\) with \(e(g)\ne0\).
Indeed after specialization the first polynomial is \(1\)
and the second is \(T^{n+1}\), so that resultant is \(1\).
If \(F(x)\notin V\) over a field, then \(a(x)=0\).
For \(x\ne0\), the two-chart relation says
\(G(x^{-1})\notin V\), so \(h(x)=0\); for \(x=0\) the last
equality holds by its factor \(T\).
Bezout's resultant identity therefore gives \(g=0\) in that
field.  Consequently, whenever \(g\) is nonzero in a field
specialization of \(R\), the whole affine-line image of \(F\)
lies in \(V\).  Interchanging \(F,G\) gives the same assertion
for \(G\).  These are the source's
\texttt{exists\_poly}, \texttt{exists\_elim}, and
\texttt{exists\_good}.

We next use its specialization at a unit.
Let \(R\) be a domain, equipped with \(L\to R\) and a surjective
\(K\)-algebra map \(e:R\to L\).
Let \(t\in R^\times\) have \(e(t)=1\).
From the two maps in \eqref{eq:faithful-torus-two-charts} form
\[
 \begin{split}
 F(X)&=b_R(tX)b_R(X)^{-1},\\
 G(X)&=b'_R(t^{-1}X)b'_R(X)^{-1}.
 \end{split}
\]
They glue under \(X\mapsto X^{-1}\) on field-valued nonzero
points.  Under \(e\) both become the identity map.
Choose an affine open \(V\) of \(B\) containing the identity.
The preceding resultant calculation for \(F,G\) and for \(G,F\)
gives two elements with nonzero specialization under \(e\).
They are nonzero in the domain \(R\), so remain nonzero in
\(Q=\operatorname{Frac}R\).
It follows that both affine-line maps over \(Q\) factor through
the same affine open \(V\).

For any regular function \(s\) on \(V\), their pullbacks are
polynomials \(P_s(X),Q_s(X)\in Q[X]\) with
\[
                          P_s(X)=Q_s(X^{-1})\quad\text{in }Q(X).
\]
This forces \(P_s\) to be constant.
For example, multiply the equality by \(X^{\deg Q_s}\);
the reflected polynomial on the right has degree at most
\(\deg Q_s\), whereas the nonzero polynomial on the left
has degree \(\deg P_s+\deg Q_s\).
The zero case is immediate.
Thus \(F_Q(1)=F_Q(0)=1\), giving
\[
                            b_R(t)=b_R(1)
\]
after passing to \(Q\).
It holds already over \(R\): for a separated target, equality
of maps from the integral affine scheme can be tested after
the injective inclusion \(R\hookrightarrow Q\).
This is \texttt{val\_unit\_eq}.

Take \(R=L[T,T^{-1}]\), \(t=T\), and \(e\) evaluation at \(T=1\).
The last equality shows that the original morphism from
\(\mathbb G_{m,L}\) is constant.  Equivalently, evaluating the
universal equality at any unit gives equality of its values at
all \(L\)-units.  This proves the one-variable assertion over
every extension \(L/K\).

Finally take a morphism \(f:D(X)\to B\).
For every extension field \(L/K\), every torus point \(z\in D(X)(L)\),
and every cocharacter \(\lambda:\mathbb G_m\to D(X)\), the map
\[
                x\longmapsto f\bigl(z\lambda(x)\bigr)
\]
is constant by the one-variable assertion.
A basis of \(X\) and its dual cocharacters express each point
as a finite product of such cocharacter values.  Successive
use of that equality gives \(f(z)=f(1)\).
To obtain equality of morphisms rather than only values, take
the universal torus point over the domain \(K[X]\), test it in
its fraction field, and use separatedness to descend the
equality along \(K[X]\hookrightarrow\operatorname{Frac}K[X]\).
This is the source's \texttt{eq\_const}.
A group homomorphism has identity value \(1\), giving the
last assertion.
\end{proof}

\paragraph{Frozen-source ledger.}
The source is
\href{https://github.com/ChristianKudeba/alpoge-rh-verification/blob/alpoge-source-7cbf7a0c/src/Development/AbelianVarietyTheory/Proofs/DegenerationProofs/SplitTorusToAbelianVarietyConst.lean}
{\texttt{SplitTorusToAbelianVarietyConst}}
(blob \texttt{ba8b522aec7afdd0c9bf4636efd4dff843172d2d}),
checked against the frozen baseline.
The argument above follows
\texttt{exists\_extension}, \texttt{valuationRing\_stalk},
\texttt{exists\_poly}, \texttt{exists\_elim},
\texttt{exists\_good}, \texttt{natDegree\_eq\_zero\_of\_inv},
\texttt{val\_eq\_of\_affine}, \texttt{val\_unit\_eq},
\texttt{exists\_extension\_line}, \texttt{ev\_eq}, and
\texttt{eq\_const}.
The eventual abelian-variety wrapper has the stronger hypotheses
of an algebraically closed field and an abelian scheme.
The proof above retains the actual stronger intermediate
proper-group statement; it does not replace its mechanism by
an Albanese or no-rational-curves assertion.
